
\documentclass{article}
\usepackage{colortbl}
\usepackage{makecell}
\usepackage{multirow}
\usepackage{supertabular}

\begin{document}

\newcounter{utterance}

\centering \large Interaction Transcript for game `function\_detective', experiment `two\_numbers\_passive\_examples', episode 295 with Qwen3.6{-}35B{-}A3B{-}FP8{-}without{-}reasoning.
\vspace{24pt}

{ \footnotesize  \setcounter{utterance}{1}
\setlength{\tabcolsep}{0pt}
\begin{supertabular}{c@{$\;$}|p{.15\linewidth}@{}p{.15\linewidth}p{.15\linewidth}p{.15\linewidth}p{.15\linewidth}p{.15\linewidth}}
    \# & \multicolumn{2}{c}{Player} && \multicolumn{2}{c}{Game Master} \\
    \hline

    \theutterance \stepcounter{utterance}  
    & & & \multicolumn{4}{p{0.6\linewidth}}{
        \cellcolor[rgb]{0.9,0.9,0.9}{
            \makecell[{{p{\linewidth}}}]{
                \texttt{\tiny{[P1$\langle$GM]}}
                \texttt{I have a hidden Python function for you to figure out.} \\
\\ 
\texttt{Your job is to infer what it does and then submit a correct Python implementation.} \\
\\ 
\texttt{Function signature:} \\
\texttt{(a: int, b: int) {-}> int} \\
\\ 
\texttt{You have 15 turns total.} \\
\\ 
\texttt{I control which valid input{-}output examples you see. You can request them one at a time.} \\
\\ 
\texttt{You must make your best SOLVE attempt no later than the second{-}to{-}last turn.} \\
\texttt{Do not wait until the very end.} \\
\\ 
\texttt{You can use these actions:} \\
\texttt{{-} JSON format described at the end of the prompt} \\
\texttt{{-} SOLVE: ```python ... ```} \\
\\ 
\texttt{Do not add import statements in your guess.} \\
\texttt{Anything you might need is already available, so just write the function body in hypothesis and return the answer.} \\
\\ 
\texttt{Preloaded examples:} \\
\texttt{1. ((1, 10), 3)} \\
\\ 
\texttt{Scoring:} \\
\texttt{{-} correctness is required for any non{-}zero score} \\
\texttt{{-} in interactive modes, solving with fewer information requests is better} \\
\\ 
\texttt{Input formatting reminders:} \\
\texttt{{-} strings must be quoted} \\
\texttt{{-} booleans must be True/False} \\
\texttt{{-} use commas between multiple values} \\
\texttt{{-} do not write any extra action tags in the same message} \\
\\ 
\texttt{Format example:} \\
\texttt{You write:} \\
\texttt{\{'current\_mode': <probing/confirming/don't know; choose whichever best describes your current state in the task>, 'input\_rationale': <brief description of why you chose this hypothesis>, 'hypothesis': <either the exact string "unknown" if you have no guess yet OR your current best guess as runnable Python matching the given signature, with a one{-}line docstring stating what it does. Never write placeholder code, if you cannot commit to a guess, use "unknown". Match the given signature.} \\
\texttt{e.g.} \\
\texttt{def f(x):} \\
\texttt{	"""Returns the absolute value of x."""} \\
\texttt{	return x if x >= 0 else {-}x} \\
\\ 
\texttt{Use "unknown" if you have no guess.>, 'confidence\_score': <score between 1{-}100. 100 means highest confidence in your hypothesis>\}} \\
\\ 
\texttt{Output only one valid json per turn.} \\
            }
        }
    }
     \\ \\

    \theutterance \stepcounter{utterance}  
    & \multicolumn{4}{p{0.6\linewidth}}{
        \cellcolor[rgb]{0.561,0.737,0.561}{
            \makecell[{{p{\linewidth}}}]{
                \texttt{\tiny{[P1$\rangle$GM]}}
                \texttt{\{'current\_mode': 'probing', 'input\_rationale': 'I have one example: f(1, 10) = 3. This is not enough to determine the function. I need more examples to understand the relationship between a, b, and the result. I will request another example with different values.', 'hypothesis': 'unknown', 'confidence\_score': 10\}} \\
            }
        }
    }
    & & \\ \\

    \theutterance \stepcounter{utterance}  
    & & & \multicolumn{4}{p{0.6\linewidth}}{
        \cellcolor[rgb]{0.9,0.9,0.9}{
            \makecell[{{p{\linewidth}}}]{
                \texttt{\tiny{[P1$\langle$GM]}}
                \texttt{Game Violation: ParseError: Response is not valid JSON.} \\
\texttt{Output ONLY:} \\
\texttt{One JSON with confidence\_score, current\_mode, input\_rationale, and hypothesis keys ...} \\
\texttt{OR} \\
\texttt{SOLVE: ```python ... ```} \\
            }
        }
    }
     \\ \\

    \theutterance \stepcounter{utterance}  
    & \multicolumn{4}{p{0.6\linewidth}}{
        \cellcolor[rgb]{0.561,0.737,0.561}{
            \makecell[{{p{\linewidth}}}]{
                \texttt{\tiny{[P1$\rangle$GM]}}
                \texttt{\{'current\_mode': 'probing', 'input\_rationale': 'I have one example: f(1, 10) = 3. This is not enough to determine the function. I need more examples to understand the relationship between a, b, and the result. I will request another example with different values.', 'hypothesis': 'unknown', 'confidence\_score': 10\}} \\
            }
        }
    }
    & & \\ \\

    \theutterance \stepcounter{utterance}  
    & & & \multicolumn{4}{p{0.6\linewidth}}{
        \cellcolor[rgb]{0.9,0.9,0.9}{
            \makecell[{{p{\linewidth}}}]{
                \texttt{\tiny{[P1$\langle$GM]}}
                \texttt{Game Violation: ParseError: Response is not valid JSON.} \\
\texttt{Output ONLY:} \\
\texttt{One JSON with confidence\_score, current\_mode, input\_rationale, and hypothesis keys ...} \\
\texttt{OR} \\
\texttt{SOLVE: ```python ... ```} \\
            }
        }
    }
     \\ \\

    \theutterance \stepcounter{utterance}  
    & \multicolumn{4}{p{0.6\linewidth}}{
        \cellcolor[rgb]{0.561,0.737,0.561}{
            \makecell[{{p{\linewidth}}}]{
                \texttt{\tiny{[P1$\rangle$GM]}}
                \texttt{\{"current\_mode": "probing", "input\_rationale": "I have one example: f(1, 10) = 3. This is not enough to determine the function. I need more examples to understand the relationship between a, b, and the result. I will request another example with different values.", "hypothesis": "unknown", "confidence\_score": 10\}} \\
            }
        }
    }
    & & \\ \\

    \theutterance \stepcounter{utterance}  
    & & & \multicolumn{4}{p{0.6\linewidth}}{
        \cellcolor[rgb]{0.9,0.9,0.9}{
            \makecell[{{p{\linewidth}}}]{
                \texttt{\tiny{[P1$\langle$GM]}}
                \texttt{OUTPUT: (({-}10, 1), 4)} \\
            }
        }
    }
     \\ \\

    \theutterance \stepcounter{utterance}  
    & \multicolumn{4}{p{0.6\linewidth}}{
        \cellcolor[rgb]{0.561,0.737,0.561}{
            \makecell[{{p{\linewidth}}}]{
                \texttt{\tiny{[P1$\rangle$GM]}}
                \texttt{\{"current\_mode": "probing", "input\_rationale": "I have two examples: f(1, 10) = 3 and f({-}10, 1) = 4. Let's analyze these.} \\
\texttt{Example 1: a=1, b=10, result=3.} \\
\texttt{Example 2: a={-}10, b=1, result=4.} \\
\\ 
\texttt{Let's look at the sum: 1+10=11, {-}10+1={-}9. No obvious link to 3 or 4.} \\
\texttt{Let's look at the difference: 10{-}1=9, 1{-}({-}10)=11. No obvious link.} \\
\texttt{Let's look at the product: 10, {-}10. No.} \\
\texttt{Let's look at absolute values: |1|+|10|=11, |{-}10|+|1|=11.} \\
\texttt{Let's look at the number of digits or something else?} \\
\texttt{Maybe it's related to the number of set bits?} \\
\texttt{bin(1)='0b1' (1 bit), bin(10)='0b1010' (2 bits). Total 3 bits. Result 3.} \\
\texttt{bin({-}10) in two's complement? Python's bin({-}10) is '{-}0b1010'. The bits are 1,0,1,0} \\
            }
        }
    }
    & & \\ \\

    \theutterance \stepcounter{utterance}  
    & & & \multicolumn{4}{p{0.6\linewidth}}{
        \cellcolor[rgb]{0.9,0.9,0.9}{
            \makecell[{{p{\linewidth}}}]{
                \texttt{\tiny{[P1$\langle$GM]}}
                \texttt{Game Violation: ParseError: Response is not valid JSON.} \\
\texttt{Output ONLY:} \\
\texttt{One JSON with confidence\_score, current\_mode, input\_rationale, and hypothesis keys ...} \\
\texttt{OR} \\
\texttt{SOLVE: ```python ... ```} \\
            }
        }
    }
     \\ \\

    \theutterance \stepcounter{utterance}  
    & \multicolumn{4}{p{0.6\linewidth}}{
        \cellcolor[rgb]{0.561,0.737,0.561}{
            \makecell[{{p{\linewidth}}}]{
                \texttt{\tiny{[P1$\rangle$GM]}}
                \texttt{\{"current\_mode": "probing", "input\_rationale": "I have two examples: f(1, 10) = 3 and f({-}10, 1) = 4. Let's analyze these.} \\
\texttt{Example 1: a=1, b=10, result=3.} \\
\texttt{Example 2: a={-}10, b=1, result=4.} \\
\\ 
\texttt{Let's look at the sum: 1+10=11, {-}10+1={-}9. No obvious link to 3 or 4.} \\
\texttt{Let's look at the difference: 10{-}1=9, 1{-}({-}10)=11. No obvious link.} \\
\texttt{Let's look at the product: 10, {-}10. No.} \\
\texttt{Let's look at absolute values: |1|+|10|=11, |{-}10|+|1|=11.} \\
\texttt{Let's look at the number of digits or something else?} \\
\texttt{Maybe it's related to the number of set bits?} \\
\texttt{bin(1)='0b1' (1 bit), bin(10)='0b1010' (2 bits). Total 3 bits. Result 3.} \\
\texttt{bin({-}10) in two's complement? Python's bin({-}10) is '{-}0b1010'. The bits are 1,0,1,0} \\
            }
        }
    }
    & & \\ \\

    \theutterance \stepcounter{utterance}  
    & & & \multicolumn{4}{p{0.6\linewidth}}{
        \cellcolor[rgb]{0.9,0.9,0.9}{
            \makecell[{{p{\linewidth}}}]{
                \texttt{\tiny{[P1$\langle$GM]}}
                \texttt{Game Violation: ParseError: Response is not valid JSON.} \\
\texttt{Output ONLY:} \\
\texttt{One JSON with confidence\_score, current\_mode, input\_rationale, and hypothesis keys ...} \\
\texttt{OR} \\
\texttt{SOLVE: ```python ... ```} \\
            }
        }
    }
     \\ \\

    \theutterance \stepcounter{utterance}  
    & \multicolumn{4}{p{0.6\linewidth}}{
        \cellcolor[rgb]{0.561,0.737,0.561}{
            \makecell[{{p{\linewidth}}}]{
                \texttt{\tiny{[P1$\rangle$GM]}}
                \texttt{\{"current\_mode": "probing", "input\_rationale": "I have two examples: f(1, 10) = 3 and f({-}10, 1) = 4. The hypothesis about set bits for positive numbers seems promising (1 has 1 bit, 10 has 2 bits, sum=3). For {-}10, if we consider the two's complement representation or just the magnitude's bits plus a sign bit, it might work. Let's test another example to confirm the bit{-}counting hypothesis, specifically checking if it's just the sum of set bits in the binary representation of the absolute values, or something involving the sign.", "hypothesis": "unknown", "confidence\_score": 30\}} \\
            }
        }
    }
    & & \\ \\

    \theutterance \stepcounter{utterance}  
    & & & \multicolumn{4}{p{0.6\linewidth}}{
        \cellcolor[rgb]{0.9,0.9,0.9}{
            \makecell[{{p{\linewidth}}}]{
                \texttt{\tiny{[P1$\langle$GM]}}
                \texttt{OUTPUT: ((10, {-}99), 4)} \\
            }
        }
    }
     \\ \\

    \theutterance \stepcounter{utterance}  
    & \multicolumn{4}{p{0.6\linewidth}}{
        \cellcolor[rgb]{0.561,0.737,0.561}{
            \makecell[{{p{\linewidth}}}]{
                \texttt{\tiny{[P1$\rangle$GM]}}
                \texttt{\{"current\_mode": "probing", "input\_rationale": "I have three examples: f(1, 10) = 3, f({-}10, 1) = 4, f(10, {-}99) = 4.} \\
\texttt{Let's check the set bit hypothesis again.} \\
\texttt{1. a=1 (0b1, 1 bit), b=10 (0b1010, 2 bits). Sum = 3. Matches.} \\
\texttt{2. a={-}10, b=1. Result 4.} \\
\texttt{   If we take absolute values: |{-}10|=10 (2 bits), |1|=1 (1 bit). Sum = 3. But result is 4.} \\
\texttt{   So it's not just sum of bits of absolute values.} \\
\texttt{   Maybe it's sum of bits of the numbers themselves in two's complement?} \\
\texttt{   In Python, integers are arbitrary precision, so two's complement is infinite.} \\
\texttt{   However, often in these puzzles, 'bits' refers to the binary representation string or a fixed width.} \\
\texttt{   Let's look at the binary strings:} \\
\texttt{   bin(1) = '0b1' {-}> 1 one.} \\
\texttt{   bin(10) = '0b1010' {-}> 2 ones.} \\
\texttt{   Total 3.} \\
\\ 
\texttt{   bin({-}10) = '{-}0b1010'. If} \\
            }
        }
    }
    & & \\ \\

    \theutterance \stepcounter{utterance}  
    & & & \multicolumn{4}{p{0.6\linewidth}}{
        \cellcolor[rgb]{0.9,0.9,0.9}{
            \makecell[{{p{\linewidth}}}]{
                \texttt{\tiny{[P1$\langle$GM]}}
                \texttt{Game Violation: ParseError: Response is not valid JSON.} \\
\texttt{Output ONLY:} \\
\texttt{One JSON with confidence\_score, current\_mode, input\_rationale, and hypothesis keys ...} \\
\texttt{OR} \\
\texttt{SOLVE: ```python ... ```} \\
            }
        }
    }
     \\ \\

    \theutterance \stepcounter{utterance}  
    & \multicolumn{4}{p{0.6\linewidth}}{
        \cellcolor[rgb]{0.561,0.737,0.561}{
            \makecell[{{p{\linewidth}}}]{
                \texttt{\tiny{[P1$\rangle$GM]}}
                \texttt{\{"current\_mode": "probing", "input\_rationale": "I have three examples: f(1, 10) = 3, f({-}10, 1) = 4, f(10, {-}99) = 4.} \\
\texttt{Hypothesis: The function returns the total number of set bits (1s) in the binary representation of the absolute values of a and b, plus 1 if either number is negative.} \\
\texttt{Check 1: a=1 (1 bit), b=10 (2 bits). Both positive. Sum = 3. Matches.} \\
\texttt{Check 2: a={-}10 (2 bits), b=1 (1 bit). a is negative. Sum = 2+1+1 = 4. Matches.} \\
\texttt{Check 3: a=10 (2 bits), b={-}99. 99 is 0b1100011 (4 bits). b is negative. Sum = 2+4+1 = 7. But result is 4.} \\
\texttt{So this hypothesis is incorrect.} \\
\\ 
\texttt{Let's re{-}examine f(10, {-}99) = 4.} \\
\texttt{a=10 (0b1010, 2 bits).} \\
\texttt{b={-}99.} \\
\texttt{Maybe it's just the number of set bits in the binary representation of the sum?} \\
\texttt{1+10=11 (0b} \\
            }
        }
    }
    & & \\ \\

    \theutterance \stepcounter{utterance}  
    & & & \multicolumn{4}{p{0.6\linewidth}}{
        \cellcolor[rgb]{0.9,0.9,0.9}{
            \makecell[{{p{\linewidth}}}]{
                \texttt{\tiny{[P1$\langle$GM]}}
                \texttt{Game Violation: ParseError: Response is not valid JSON.} \\
\texttt{Output ONLY:} \\
\texttt{One JSON with confidence\_score, current\_mode, input\_rationale, and hypothesis keys ...} \\
\texttt{OR} \\
\texttt{SOLVE: ```python ... ```} \\
            }
        }
    }
     \\ \\

    \theutterance \stepcounter{utterance}  
    & \multicolumn{4}{p{0.6\linewidth}}{
        \cellcolor[rgb]{0.561,0.737,0.561}{
            \makecell[{{p{\linewidth}}}]{
                \texttt{\tiny{[P1$\rangle$GM]}}
                \texttt{\{"current\_mode": "probing", "input\_rationale": "I have three examples: f(1, 10) = 3, f({-}10, 1) = 4, f(10, {-}99) = 4.} \\
\texttt{Let's check the sum of set bits in the binary representation of the absolute values again, but maybe I miscalculated 99.} \\
\texttt{99 in binary: 64 + 32 + 2 + 1 = 0b1100011. That has 4 bits set.} \\
\texttt{So for (10, {-}99): |10| has 2 bits, |{-}99| has 4 bits. Sum = 6. Result is 4.} \\
\texttt{This doesn't match simple bit counting.} \\
\\ 
\texttt{Let's look at the difference between the numbers.} \\
\texttt{1. 1, 10 {-}> diff 9. Result 3.} \\
\texttt{2. {-}10, 1 {-}> diff 11. Result 4.} \\
\texttt{3. 10, {-}99 {-}> diff 109. Result 4.} \\
\\ 
\texttt{Let's look at the product.} \\
\texttt{1. 10. Result 3.} \\
\texttt{2. {-}10. Result 4.} \\
\texttt{3. {-}990. Result 4.} \\
\\ 
\texttt{Let's look at the number of digits in decimal?} \\
\texttt{1. 1 (} \\
            }
        }
    }
    & & \\ \\

    \theutterance \stepcounter{utterance}  
    & & & \multicolumn{4}{p{0.6\linewidth}}{
        \cellcolor[rgb]{0.9,0.9,0.9}{
            \makecell[{{p{\linewidth}}}]{
                \texttt{\tiny{[P1$\langle$GM]}}
                \texttt{Game Violation: ParseError: Response is not valid JSON.} \\
\texttt{Output ONLY:} \\
\texttt{One JSON with confidence\_score, current\_mode, input\_rationale, and hypothesis keys ...} \\
\texttt{OR} \\
\texttt{SOLVE: ```python ... ```} \\
            }
        }
    }
     \\ \\

    \theutterance \stepcounter{utterance}  
    & \multicolumn{4}{p{0.6\linewidth}}{
        \cellcolor[rgb]{0.561,0.737,0.561}{
            \makecell[{{p{\linewidth}}}]{
                \texttt{\tiny{[P1$\rangle$GM]}}
                \texttt{\{"current\_mode": "probing", "input\_rationale": "I have three examples: f(1, 10) = 3, f({-}10, 1) = 4, f(10, {-}99) = 4.} \\
\texttt{Let's try a different hypothesis: The function returns the number of set bits in the binary representation of the sum of the absolute values of a and b.} \\
\texttt{1. |1| + |10| = 11. bin(11) = '0b1011'. Set bits = 3. Matches.} \\
\texttt{2. |{-}10| + |1| = 11. bin(11) = '0b1011'. Set bits = 3. But result is 4. Mismatch.} \\
\\ 
\texttt{Let's try: The function returns the number of set bits in the binary representation of the difference between the absolute values?} \\
\texttt{1. |10| {-} |1| = 9. bin(9) = '0b1001'. Set bits = 2. Mismatch.} \\
\\ 
\texttt{Let's try: The function returns the number of set bits in the binary representation of the product of the absolute values?} \\
\texttt{1. 1 * 10 = 10. bin(10) = '0b1010'. Set bits = 2. Mismatch.} \\
\\ 
\texttt{Let's try:} \\
            }
        }
    }
    & & \\ \\

    \theutterance \stepcounter{utterance}  
    & & & \multicolumn{4}{p{0.6\linewidth}}{
        \cellcolor[rgb]{0.9,0.9,0.9}{
            \makecell[{{p{\linewidth}}}]{
                \texttt{\tiny{[P1$\langle$GM]}}
                \texttt{Game Violation: ParseError: Response is not valid JSON.} \\
\texttt{Output ONLY:} \\
\texttt{One JSON with confidence\_score, current\_mode, input\_rationale, and hypothesis keys ...} \\
\texttt{OR} \\
\texttt{SOLVE: ```python ... ```} \\
            }
        }
    }
     \\ \\

    \theutterance \stepcounter{utterance}  
    & \multicolumn{4}{p{0.6\linewidth}}{
        \cellcolor[rgb]{0.561,0.737,0.561}{
            \makecell[{{p{\linewidth}}}]{
                \texttt{\tiny{[P1$\rangle$GM]}}
                \texttt{\{"current\_mode": "probing", "input\_rationale": "I have three examples: f(1, 10) = 3, f({-}10, 1) = 4, f(10, {-}99) = 4.} \\
\texttt{Let's look at the binary representations of the numbers themselves, including the sign bit if negative, but treating them as standard integers.} \\
\texttt{Maybe it's the number of set bits in the XOR of a and b?} \\
\texttt{1. 1 \^{} 10 = 11 (0b1011). Set bits = 3. Matches.} \\
\texttt{2. {-}10 \^{} 1. In Python, {-}10 is ...11110110 (infinite 1s to the left). 1 is ...00000001.} \\
\texttt{   XOR: ...11110111. This has infinite set bits. This doesn't seem right for a simple puzzle unless we consider a fixed width or just the magnitude.} \\
\\ 
\texttt{Let's reconsider the set bits of absolute values.} \\
\texttt{1. |1|=1 (1 bit), |10|=10 (2 bits). Sum=3.} \\
\texttt{2. |{-}10|=10 (2 bits), |1|=1 (1 bit). Sum=3. Result=4.} \\
\texttt{3. |10|=10 (2 bits), |{-}99} \\
            }
        }
    }
    & & \\ \\

    \theutterance \stepcounter{utterance}  
    & & & \multicolumn{4}{p{0.6\linewidth}}{
        \cellcolor[rgb]{0.9,0.9,0.9}{
            \makecell[{{p{\linewidth}}}]{
                \texttt{\tiny{[P1$\langle$GM]}}
                \texttt{Game Violation: ParseError: Response is not valid JSON.} \\
\texttt{Output ONLY:} \\
\texttt{One JSON with confidence\_score, current\_mode, input\_rationale, and hypothesis keys ...} \\
\texttt{OR} \\
\texttt{SOLVE: ```python ... ```} \\
            }
        }
    }
     \\ \\

    \theutterance \stepcounter{utterance}  
    & \multicolumn{4}{p{0.6\linewidth}}{
        \cellcolor[rgb]{0.561,0.737,0.561}{
            \makecell[{{p{\linewidth}}}]{
                \texttt{\tiny{[P1$\rangle$GM]}}
                \texttt{\{"current\_mode": "probing", "input\_rationale": "I have three examples: f(1, 10) = 3, f({-}10, 1) = 4, f(10, {-}99) = 4.} \\
\texttt{Let's look at the number of set bits in the binary representation of the sum of the numbers.} \\
\texttt{1. 1 + 10 = 11. bin(11) = '0b1011'. Set bits = 3. Matches.} \\
\texttt{2. {-}10 + 1 = {-}9. bin({-}9) is not standard for bit counting.} \\
\texttt{   However, if we look at the absolute value of the sum: |{-}9| = 9. bin(9) = '0b1001'. Set bits = 2. Mismatch.} \\
\\ 
\texttt{Let's look at the number of set bits in the binary representation of the difference.} \\
\texttt{1. 10 {-} 1 = 9. bin(9) = '0b1001'. Set bits = 2. Mismatch.} \\
\texttt{2. 1 {-} ({-}10) = 11. bin(11) = '0b1011'. Set bits = 3. Mismatch.} \\
\\ 
\texttt{Let's look at the number of set bits in the binary representation of the product.} \\
\texttt{1. 1 * 10} \\
            }
        }
    }
    & & \\ \\

    \theutterance \stepcounter{utterance}  
    & & & \multicolumn{4}{p{0.6\linewidth}}{
        \cellcolor[rgb]{0.9,0.9,0.9}{
            \makecell[{{p{\linewidth}}}]{
                \texttt{\tiny{[P1$\langle$GM]}}
                \texttt{Game Violation: ParseError: Response is not valid JSON.} \\
\texttt{Output ONLY:} \\
\texttt{One JSON with confidence\_score, current\_mode, input\_rationale, and hypothesis keys ...} \\
\texttt{OR} \\
\texttt{SOLVE: ```python ... ```} \\
            }
        }
    }
     \\ \\

    \theutterance \stepcounter{utterance}  
    & \multicolumn{4}{p{0.6\linewidth}}{
        \cellcolor[rgb]{0.561,0.737,0.561}{
            \makecell[{{p{\linewidth}}}]{
                \texttt{\tiny{[P1$\rangle$GM]}}
                \texttt{\{"current\_mode": "probing", "input\_rationale": "I have three examples: f(1, 10) = 3, f({-}10, 1) = 4, f(10, {-}99) = 4.} \\
\texttt{Let's try one more hypothesis: The function returns the number of set bits in the binary representation of the sum of the absolute values of a and b, but with a twist for negative numbers.} \\
\texttt{Actually, let's look at the number of set bits in the binary representation of the XOR of the absolute values.} \\
\texttt{1. |1| \^{} |10| = 1 \^{} 10 = 11 (0b1011). Set bits = 3. Matches.} \\
\texttt{2. |{-}10| \^{} |1| = 10 \^{} 1 = 11 (0b1011). Set bits = 3. But result is 4. Mismatch.} \\
\\ 
\texttt{Let's look at the number of set bits in the binary representation of the sum of the absolute values again.} \\
\texttt{1. |1| + |10| = 11. Set bits = 3. Matches.} \\
\texttt{2. |{-}10| + |1| = 11. Set bits = 3. Mismatch.} \\
\\ 
\texttt{Let's look at the number of set bits in the binary representation of the difference of the absolute values.} \\
\texttt{1.} \\
            }
        }
    }
    & & \\ \\

    \theutterance \stepcounter{utterance}  
    & & & \multicolumn{4}{p{0.6\linewidth}}{
        \cellcolor[rgb]{0.9,0.9,0.9}{
            \makecell[{{p{\linewidth}}}]{
                \texttt{\tiny{[P1$\langle$GM]}}
                \texttt{Game Violation: ParseError: Response is not valid JSON.} \\
\texttt{Output ONLY:} \\
\texttt{One JSON with confidence\_score, current\_mode, input\_rationale, and hypothesis keys ...} \\
\texttt{OR} \\
\texttt{SOLVE: ```python ... ```} \\
            }
        }
    }
     \\ \\

    \theutterance \stepcounter{utterance}  
    & \multicolumn{4}{p{0.6\linewidth}}{
        \cellcolor[rgb]{0.561,0.737,0.561}{
            \makecell[{{p{\linewidth}}}]{
                \texttt{\tiny{[P1$\rangle$GM]}}
                \texttt{\{"current\_mode": "probing", "input\_rationale": "I have three examples: f(1, 10) = 3, f({-}10, 1) = 4, f(10, {-}99) = 4.} \\
\texttt{Let's try a different approach. Maybe it's related to the number of digits in the decimal representation?} \\
\texttt{1. 1 (1 digit), 10 (2 digits). Sum = 3. Matches.} \\
\texttt{2. {-}10 (3 chars, but 2 digits), 1 (1 digit). If we count digits: 2 + 1 = 3. Result 4. Mismatch.} \\
\texttt{   If we count characters: 3 + 1 = 4. Matches.} \\
\texttt{3. 10 (2 digits), {-}99 (3 chars, 2 digits).} \\
\texttt{   Digits: 2 + 2 = 4. Matches.} \\
\texttt{   Characters: 2 + 3 = 5. Mismatch.} \\
\\ 
\texttt{So, hypothesis: The function returns the sum of the number of digits in the decimal representation of the absolute values of a and b.} \\
\texttt{1. |1| has 1 digit. |10| has 2 digits. Sum = 3. Matches.} \\
\texttt{2. |{-}10| has 2 digits. |1| has 1 digit. Sum = 3. Mismatch (Result is} \\
            }
        }
    }
    & & \\ \\

    \theutterance \stepcounter{utterance}  
    & & & \multicolumn{4}{p{0.6\linewidth}}{
        \cellcolor[rgb]{0.9,0.9,0.9}{
            \makecell[{{p{\linewidth}}}]{
                \texttt{\tiny{[P1$\langle$GM]}}
                \texttt{Game Violation: ParseError: Response is not valid JSON.} \\
\texttt{Output ONLY:} \\
\texttt{One JSON with confidence\_score, current\_mode, input\_rationale, and hypothesis keys ...} \\
\texttt{OR} \\
\texttt{SOLVE: ```python ... ```} \\
            }
        }
    }
     \\ \\

    \theutterance \stepcounter{utterance}  
    & \multicolumn{4}{p{0.6\linewidth}}{
        \cellcolor[rgb]{0.561,0.737,0.561}{
            \makecell[{{p{\linewidth}}}]{
                \texttt{\tiny{[P1$\rangle$GM]}}
                \texttt{\{"current\_mode": "probing", "input\_rationale": "I have three examples: f(1, 10) = 3, f({-}10, 1) = 4, f(10, {-}99) = 4.} \\
\texttt{Let's re{-}evaluate the digit count hypothesis.} \\
\texttt{1. |1| has 1 digit. |10| has 2 digits. Sum = 3. Matches.} \\
\texttt{2. |{-}10| has 2 digits. |1| has 1 digit. Sum = 3. Mismatch (Result is 4).} \\
\texttt{3. |10| has 2 digits. |{-}99| has 2 digits. Sum = 4. Matches.} \\
\\ 
\texttt{The second example is the outlier for simple digit counting.} \\
\texttt{Let's look at the binary representation again.} \\
\texttt{1. 1 (0b1), 10 (0b1010).} \\
\texttt{2. {-}10, 1.} \\
\texttt{3. 10, {-}99.} \\
\\ 
\texttt{Maybe it's the number of set bits in the binary representation of the sum of the numbers, but using two's complement for a fixed width? Or maybe just the number of set bits in the binary representation of the absolute value of the sum?} \\
\texttt{1. 1+10=11. bin(11)='0b1011'. 3 bits. Matches.} \\
            }
        }
    }
    & & \\ \\

    \theutterance \stepcounter{utterance}  
    & & & \multicolumn{2}{p{0.3\linewidth}}{
        \cellcolor[rgb]{0.9,0.9,0.9}{
            \makecell[{{p{\linewidth}}}]{
                \texttt{\tiny{[GM$|$GM]}}
                \texttt{Game Verdict: LOSS (turn limit reached)} \\
            }
        }
    }
    & & \\ \\

\end{supertabular}
}

\end{document}
